% ============================================================
%  Section 4: Results
% ============================================================

\section{Results}

\begin{frame}{Result 1: Generation Engines Do Not Automatically Win}
  \small
  \begin{tabularx}{\textwidth}{llYYY}
    \toprule
    \textbf{Engine} & \textbf{Variant} & \textbf{MB/s} & \textbf{Ratio} & \textbf{Speedup} \\
    \midrule
    Transformers & + KV cache & \metric{0.0344} & 7.0492 & $1.00\times$ \\
    Transformers & No KV cache & 0.0030 & 7.0502 & $0.09\times$ \\
    vLLM & Default & 0.0130 & 7.0304 & $0.37\times$ \\
    TensorRT-LLM & Batch 256 & 0.0081 & 7.0499 & $0.023\times$ \\
    TensorRT-LLM & Batch 512 & 0.0065 & 7.9370 & $0.019\times$ \\
    \bottomrule
  \end{tabularx}

  \vspace{0.8em}
  \begin{block}{Interpretation}
    vLLM and TensorRT-LLM are excellent generation systems, but this workload exposes scheduling, synchronization, and score-capture costs that ordinary serving can hide \parencite{kwonEfficientMemoryManagement2023,Kipf2026WaitingTD}.
  \end{block}
\end{frame}

\begin{frame}{The Bottleneck Moves}
  \centering
  \begin{tikzpicture}[
    node distance=1.15cm,
    stage/.style={draw, rounded corners=2pt, minimum width=3.25cm,
                  minimum height=0.85cm, align=center, font=\small},
    arr/.style={-{Stealth[length=5pt]}, thick},
  ]
    \node[stage, fill=blue!12] (low) {Low batch\\inference latency};
    \node[stage, fill=orange!16, right=of low] (mid) {Larger batch\\AC preprocessing};
    \node[stage, fill=green!14, right=of mid] (high) {After O2/O3\\score extraction / inference};

    \draw[arr] (low) -- node[above, font=\tiny] {batching} (mid);
    \draw[arr] (mid) -- node[above, font=\tiny] {optimized coding} (high);
  \end{tikzpicture}

  \vspace{1.2em}
  \begin{itemize}
    \item Larger batches make GPU work more efficient, but expose arithmetic-coding cost.
    \item O2 reduces Python and range-coder overhead.
    \item O3 removes most probability-vector processing and transfer.
  \end{itemize}
\end{frame}

\begin{frame}{Result 2: Target-Interval Capture Dominates}
  \small
  \begin{tabularx}{\textwidth}{llYYY}
    \toprule
    \textbf{Engine} & \textbf{Configuration} & \textbf{MB/s} & \textbf{Rel. speedup} & \textbf{AC time} \\
    \midrule
    Transformers & Baseline AC & 0.0018 & $1.00\times$ & 97.1\% \\
    Transformers & + O3 & \metric{0.0347} & \metric{$18.96\times$} & \metric{1.0\%} \\
    Transformers & + O2 + O3 & 0.0345 & $18.84\times$ & 1.0\% \\
    \midrule
    vLLM & Baseline AC & 0.0031 & $1.00\times$ & 89.9\% \\
    vLLM & + O1 & \metric{0.0116} & \metric{$3.43\times$} & 9.2\% \\
    vLLM & + O3 & 0.0065 & $1.94\times$ & 3.0\% \\
    \midrule
    TensorRT-LLM & Baseline AC & 0.0020 & $1.00\times$ & 58.2\% \\
    TensorRT-LLM & + O3 & \metric{0.0044} & \metric{$2.18\times$} & 0.7\% \\
    \bottomrule
  \end{tabularx}
\end{frame}

\begin{frame}{Result 3: vLLM Needs Windowing}
  \begin{columns}[T, onlytextwidth]
    \begin{column}{0.56\textwidth}
      \includegraphics[width=\textwidth]{\plotsdir/vllm_window_sweep_heatmap.pdf}
    \end{column}
    \hfill
    \begin{column}{0.40\textwidth}
      \textbf{Takeaway}
      \begin{itemize}
        \item One next-token request per step pays too much fixed overhead.
        \item Teacher forcing collects several score vectors in one call.
        \item The benefit is largest at small-to-medium batch sizes.
        \item Window size $w=8$ is a practical operating point.
      \end{itemize}
    \end{column}
  \end{columns}
\end{frame}

\begin{frame}{Result 4: Context Length Mostly Buys Ratio}
  \small
  \begin{tabularx}{0.78\textwidth}{YYYY}
    \toprule
    \textbf{Context} & \textbf{Retain} & \textbf{MB/s} & \textbf{Ratio} \\
    \midrule
    32 & 16 & 0.0324 & 6.2288 \\
    64 & 32 & 0.0338 & 6.6592 \\
    128 & 64 & 0.0339 & 6.9239 \\
    \textbf{256} & \textbf{100} & \textbf{0.0337} & \textbf{7.0600} \\
    512 & 256 & 0.0349 & 7.1007 \\
    \bottomrule
  \end{tabularx}

  \vspace{0.8em}
  \begin{itemize}
    \item Throughput is nearly flat from context 64 to 512.
    \item Compression ratio improves, with diminishing returns after the first few hundred tokens.
    \item Context 256 is a memory-efficient default for this setup.
  \end{itemize}
\end{frame}

\begin{frame}{Result 5: LLM Ratios Are Strong, But Throughput Is Not}
  \small
  \begin{tabularx}{\textwidth}{lYYY}
    \toprule
    \textbf{Codec / encoding} & \textbf{MB/s} & \textbf{Ratio} & \textbf{Notes} \\
    \midrule
    Target interval AC & 0.0349 & \metric{7.1007} & LLM, optimized AC \\
    Bitpacked ranks & 0.0363 & 2.8735 & Simpler, much worse ratio \\
    Huffman ranks & 0.0359 & 6.1962 & No throughput advantage \\
    \midrule
    LZ4 & \metric{666.66} & 1.83 & CPU baseline \\
    zstd & \metric{586.70} & 3.06 & CPU baseline \\
    FSST & 294.11 & 2.08 & DB-centric baseline \\
    xz -9e & 1.34 & 4.30 & Stronger, slower baseline \\
    \bottomrule
  \end{tabularx}

  \vspace{0.5em}
  \small LLM compression currently trades orders of magnitude of throughput for substantially higher compression ratio. Classical baselines include zstd and FSST \parencite{zstd,fsst}.
\end{frame}

\begin{frame}{Dataset Sensitivity: Code Compresses Especially Well}
  \begin{columns}[T, onlytextwidth]
    \begin{column}{0.50\textwidth}
      \small
      \begin{tabularx}{\textwidth}{lYY}
        \toprule
        \textbf{Dataset} & \textbf{MB/s} & \textbf{Ratio} \\
        \midrule
        \texttt{text8} & 0.005436 & 7.739 \\
        \texttt{code} & 0.004132 & \metric{11.203} \\
        \texttt{Yelp} & 0.002490 & 7.028 \\
        \bottomrule
      \end{tabularx}
    \end{column}
    \hfill
    \begin{column}{0.44\textwidth}
      \begin{block}{Why this matters}
        Domain-structured text can be much more predictable than generic natural language.
      \end{block}
      \begin{itemize}
        \item Logs, code, schemas, and repeated business records are plausible targets.
        \item Fine-tuned or specialized models may shift the economics further \parencite{finezip}.
      \end{itemize}
    \end{column}
  \end{columns}
\end{frame}
